% CV template formatted to match the reference image
\documentclass[11pt,letterpaper]{article}

\usepackage[margin=0.62in,top=0.38in,bottom=0.55in]{geometry}
\usepackage[T1]{fontenc}
\usepackage{lmodern}
\usepackage{microtype}
\usepackage{xcolor}
\usepackage{hyperref}
\usepackage{tabularx}
\usepackage{array}
\usepackage{enumitem}
\usepackage{multicol}

\definecolor{cvcolor}{RGB}{90, 60, 35}
\definecolor{rulecolor}{RGB}{125, 85, 50}

\hypersetup{colorlinks=true,urlcolor=cvcolor,linkcolor=cvcolor}
\pagestyle{empty}
\setlength{\parindent}{0pt}
\setlength{\parskip}{0pt}
\renewcommand{\baselinestretch}{1.02}

\setlist[itemize]{
  leftmargin=1.2em,
  topsep=2pt,
  itemsep=1pt,
  parsep=0pt
}

\newcommand{\CVHeader}{%
  \begin{minipage}[t]{0.62\textwidth}
    {\fontsize{28}{34}\selectfont\color{cvcolor} Hassan \textsc{AL KHAWALDEH}}\\[1.55em]
    {\fontsize{15}{18}\selectfont\itshape\color{cvcolor} Curriculum Vitae}
  \end{minipage}%
  \hfill
  \hspace{-2mm}
  \begin{minipage}[t]{0.34\textwidth}
    \raggedleft\small\itshape\color{cvcolor}
    Mechanical Engineering \& Mechanics\\
    Lehigh University, Bethlehem, USA\\[-1pt]
    \href{mailto:alkhawaldeh@lehigh.edu}{alkhawaldeh@lehigh.edu}\\[2pt]
     Linkedin: \href{https://www.linkedin.com/in/alkhawaldehh/}{alkhawaldeh}
  \end{minipage}
  \vspace{3.9em}
}

\newcommand{\CVSection}[1]{%
  \vspace{0.4em}
  \noindent
  \begin{tabularx}{\textwidth}{@{}p{0.82in}@{\hspace{0.14in}}X@{}}
    \raisebox{-1.45ex}{\color{rulecolor}\rule{0.82in}{4.2pt}} &
    {\fontsize{14}{16}\selectfont\color{cvcolor}#1}\\
  \end{tabularx}
  \vspace{0.55em}
}

\newcommand{\CVSubsection}[1]{%
  \vspace{0.25em}
  {\large\itshape\color{cvcolor}#1}\\[0.25em]
}

\newcommand{\EducationEntry}[6]{%
  \begin{tabularx}{\textwidth}{@{}p{0.83in}@{\hspace{0.15in}}X@{}}
    \raggedleft #1 &
    \textbf{#2}, #3.\\[-1pt]
    & Dissertation title: {\itshape #4}.\\[-1pt]
    & Advisor: #5.\\[-1pt]
  \end{tabularx}
}

\newcommand{\UndergradEntry}[6]{%
  \begin{tabularx}{\textwidth}{@{}p{0.83in}@{\hspace{0.15in}}X@{}}
    \raggedleft #1 &
    \textbf{#2}, #3.\\[-1pt]
    & Minors: #4.\\[-1pt]
  \end{tabularx}
}

\newcommand{\CVEntry}[5]{%
  \begin{tabularx}{\textwidth}{@{}p{0.98in}@{\hspace{0.12in}}X@{}}
    \raggedleft #1 &
    \textbf{#2}, #3%
    \if\relax\detokenize{#4}\relax\else, #4\fi.\\[-1pt]
    & #5\\[5pt]
  \end{tabularx}
}

\newcommand{\CVSubtitle}[1]{%
  \noindent\hspace{0.96in}%
  \raisebox{-1.0em}{\itshape\color{cvcolor}#1}\\[0.45em]
}

\newcommand{\CVItem}[2]{%
  \begin{tabularx}{\textwidth}{@{}p{0.83in}@{\hspace{0.15in}}X@{}}
    \raggedleft{#1} & #2\\[3pt]
  \end{tabularx}
}

\newcommand{\CVDoubleItem}[4]{%
  \begin{tabularx}{\textwidth}{@{}p{0.18\textwidth}@{\hspace{0.04\textwidth}}X
                                  p{0.18\textwidth}@{\hspace{0.04\textwidth}}X@{}}
    \textit{#1} & #2 & \textit{#3} & #4\\[3pt]
  \end{tabularx}
}

\begin{document}

\CVHeader
\vspace{-10mm}

% ---------------------------------------------------------
\CVSection{Education}

\EducationEntry
  {\hspace{-5mm}\mbox{2021--Ongoing}}
  {Ph.D. in Mechanical Engineering}
  {\textit{Lehigh University}, Bethlehem, PA, USA}
  {Physics-Based and Data-Driven Predictive Modeling and Control of Plasma Density in Tokamaks}
  {Dr. Eugenio Schuster}

\UndergradEntry
  {2017--2021}
  {B.S. in Mechanical Engineering}
  {\textit{Lehigh University}, Bethlehem, PA, USA}
  {Energy Engineering and Business Management}



% ---------------------------------------------------------
\CVSection{Research Interest}

\CVItem{Topics}{  Adaptive Control (model-reference, self-tuning);
  Optimal Control (LQR, MPC);
  Robust Control ($H_{\infty}$, structured uncertainty);
  Nonlinear Control (feedback linearization, Lyapunov-based methods);
  % System Theory & Identification
  System Identification (Koopman operator theory, DMD, subspace methods, gray-box modeling);
  % Learning & Computation
  Reinforcement Learning (Model-free, safe RL);
  Neural-Network Modeling (physics-informed, black-box).}
  
% ---------------------------------------------------------
\CVSection{Publications: Journal Articles}

\CVItem{[1]} {H. Al Khawaldeh, A. Pajares, S. T. Paruchuri, V. Graber, T. Rafiq, E. Schuster, J.-W. Juhn, T. Pederson, D. Shiraki, F. Turco, ``Experimental Demonstration of Reactor-Relevant  Adaptive Density Control on the DIII-D Tokamak with Coordinated Gas and Pellet Fueling,'' submitted to \textit{Plasma Physics and Controlled Fusion}, 2026. }

\CVItem{[2]} {S. T. Paruchuri, I. Ward, N. Rist, V. Graber, H. Al Khawaldeh, Z. Wang, E. Schuster,
A. Pajares, and J.-W. Juhn, ``Density Regulation With Disruption Avoidance in Next-Generation
Tokamaks Using a Safe Reinforcement Learning-Based Controller,'' \textit{Fusion Engineering and
Design}, vol. 216, 2025.}

\CVItem{[3]} {J. W. Berkery, (H. Al Khawaldeh), et al. (Collaboration Paper), ``NSTX-U Research Advancing the Physics of Spherical Tokamak,'' \textit{Nuclear Fusion}, vol. 64, 2024.}

\CVItem{[4]} {S. T. Paruchuri, V. Graber, H. Al Khawaldeh, E. Schuster, ``Dynamic Actuator Allocation via Reinforcement Learning for Concurrent Plasma Control Objectives,'' \textit{IEEE Transactions on Plasma Science}, vol. 52, 2024.}

\CVItem{[5]} {H. Al Khawaldeh, B. Leard, S.T. Paruchuri, T. Rafiq and E. Schuster, ``Model-Based Linear-Quadratic-Integral Controller for Simultaneous Regulation of the Current Profile and Normalized Beta in NSTX-U,'' \textit{Fusion Engineering and Design}, vol. 192, 2023. }

% ---------------------------------------------------------
\CVSection{Publications: Refereed Conference Abstracts}

\CVItem{[1]}{H. Al Khawaldeh, S. T. Paruchuri, V. Graber, T. Rafiq, E. Schuster, A. Pajares, J. W. Juhn, S. H. Park, ``Demonstration of Real-Time Multi-Actuator Plasma Density Control in KSTAR via Adaptive Fueling,'' \textit{34th Symposium on Fusion Technology (SOFT)}, Aix-en-Provence, France, September 21-25, 2026.}

\CVItem{[2]}{H. Al Khawaldeh, V. Graber, S. T. Paruchuri, T. Rafiq, E. Schuster, ``Coordinated Gas and Pellet Fueling for Plasma Density in KSTAR Under Reactor-Relevant Conditions,'' \textit{67th Division of Plasma Physics (DPP) Annual Meeting of the American Physical Society (APS)}, Long Beach, CA, USA, November 17-21, 2025.} 

\CVItem{[3]}{A. Pajares, H. Al Khawaldeh, F. Turco, D. Shiraki, T. Pederson, D. Eldon, J-W. Juhn, V. Graber, S. T. Paruchuri, E. Schuster, ``ITER-Relevant Density Control with Pellets in DIII-D’s ITER Baseline Scenario,'' \textit{67th Division of Plasma Physics Annual (DPP) Meeting of the American Physical Society (APS)}, Long Beach, CA, USA, November 17-21, 2025.}

\CVItem{[4]}{S. T. Paruchuri, H. Al Khawaldeh, V. Graber, E. Schuster, ``Nonlinear Adaptive Control of Deuterium-Tritium Density Ratio in ITER,'' \textit{67th Division of Plasma Physics (DPP) Annual Meeting of the American Physical Society (APS)}, Long Beach, CA, USA, November 17-21, 2025.}

\CVItem{[5]}{Y. Tao, B. Leard, H. Al Khawaldeh, T. Rafiq, E. Schuster, ``Predictive Transport Modeling of KSTAR Advanced Scenarios Using COTSIM,'' \textit{67th Division of Plasma Physics (DPP) Annual Meeting of the American Physical Society (APS)}, Long Beach, CA, USA, November 17-21, 2025.}

\CVItem{[6]}{Z. Wang, S. T. Paruchuri, H. Al Khawaldeh, T. Rafiq, E. Schuster, ``Pellet Injection-Based Model Predictive Control of the Density Profile in Tokamaks by Leveraging Deep Reinforcement Learning,'' \textit{31st IEEE Symposium on Fusion Energy (SOFE)}, Cambridge, MA, USA, June 23-26, 2025.}

\CVItem{[7]}{H. Al Khawaldeh, S.T. Paruchuri, E. Schuster, A. Pajares, and J. W. Juhn, ``Adaptive Control of Plasma Density in Tokamaks Using Coordinated Gas Puffing and Pellet Injection,'' \textit{31st IEEE Symposium on Fusion Energy (SOFE)}, Cambridge, MA, USA, June 23- 26, 2025.}

\CVItem{[8]}{H. Al Khawaldeh, S.T. Paruchuri, V. Graber, T. Rafiq, E. Schuster, A. Pajares, and J. W. Juhn, ``Advanced Density Regulation for ITER Emulated Scenarios on DIII-D via Adaptive Control Techniques,'' \textit{66th Division of Plasma Physics Annual (DPP) Meeting of the American Physical Society (APS)}, Atlanta, GA, USA, October 7-11, 2024.}

\CVItem{[9]}{S. T. Paruchuri, H. Al Khawaldeh, V. Graber, Z. Wang, N. Rist, I. Ward, T. Rafiq, E. Schuster, A. Pajares, J.W. Juhn, ``Assessing Model-based and Learning-based Strategies for Active Density Regulation in ITER,'' \textit{66th Division of Plasma Physics (DPP) Annual Meeting of the American Physical Society (APS)}, Atlanta, GA, USA, October 7-11, 2024.}

\CVItem{[10]}{H. Al Khawaldeh, S.T. Paruchuri, T. Rafiq and E. Schuster, ``Simultaneous Regulation of the Dimensionless Gain and Density by Real-time Estimation of Confinement Times,'' \textit{33nd Symposium on Fusion Technology (SOFT)}, Dublin, Ireland, September 22-27, 2024.}

\CVItem{[11]}{S. T. Paruchuri, H. Al Khawaldeh, V. Graber, E. Schuster, A. Pajares, J-W. Juhn, ``Safe Reinforcement Learning-based Controller for Disruption Free Regulation of Plasma Density in Next-generation Tokamaks,'' \textit{33rd Symposium on Fusion Technology (SOFT)}, Dublin, Ireland, September 22-27, 2024. }

\CVItem{[12]}{E. Schuster, S. T. Paruchuri, H. Al Khawaldeh, V. Graber, T. Sing, A. Pajares, J-W. Juhn, S. Park, ``Advanced Model-based Density Control on KSTAR as Testbed for ITER Operation,'' \textit{3rd International Fusion and Plasma Conference (iFPC)}, Seoul, South Korea, June 24-28, 2024. }

\CVItem{[13]}{H. Al Khawaldeh, S.T. Paruchuri, T. Rafiq and E. Schuster, ``Model-Based Control of the Dimensionless Gain in KSTAR by Leveraging Real-time Estimation of the Confinement Time,'' \textit{65th Division of Plasma Physics (DPP) Annual Meeting of the American Physical Society (APS)}, Denver, CO, USA, October 30-November 3, 2023. }

\CVItem{[14]}{H. Al Khawaldeh, S.T. Paruchuri, Z. Wang, T. Rafiq, E. Schuster, ``Simultaneous Optimal Regulation of Kinetic+Magnetic Scalar Plasma Properties for Robust Sustainment of Advanced Scenarios in NSTX-U,'' \textit{29th IAEA Fusion Energy Conference}, London, UK, October 16-21, 2023. }

\CVItem{[15]}{L. Yang, H. Al Khawaldeh, S.T. Paruchuri, X. Song, Z. Wang, E. Schuster, ``Towards Density Profile Regulation via Pellet Injection in Tokamaks Following a Hybrid Model Predictive Control Approach,'' \textit{29th IAEA Fusion Energy Conference}, London, UK, October 16-21, 2023.}

\CVItem{[16]}{J. W. Berkery, (H. Al Khawaldeh), et al. (Collaboration Paper), ``NSTX-U Research Advancing the Physics of Spherical Tokamaks,'' \textit{29th IAEA Fusion Energy Conference}, London, UK, October 16-21, 2023.}

\CVItem{[17]}{H. Al Khawaldeh, B. Leard, S.T. Paruchuri, T. Rafiq and E. Schuster, ``Model-Based Optimal Control of Core Kinetic+Magnetic Profiles and Scalar Plasma Properties in NSTX-U,'' \textit{64th Division of Plasma Physics (DPP) Annual Meeting of the American Physical Society (APS)}, Spokane, WA, USA, October 17-21, 2022.}

\CVItem{[18]}{B. Leard, C. Clauser, H. Al Khawaldeh, T. Rafiq, E. Schuster, ``Development of COTSIM+NUBEAMNet One-Dimensional Predictive Simulation Capability for Scenario Optimization in NSTX-U,'' \textit{64th Division of Plasma Physics (DPP) Annual Meeting of the American Physical Society (APS)}, Spokane, WA, USA, October 17-21, 2022.}

\CVItem{[19]}{H. Al Khawaldeh, B. Leard, S.T. Paruchuri, T. Rafiq and E. Schuster, ``Model-Based Linear-Quadratic-Integral Controller for Simultaneous Regulation of Current Profile and Normalized Beta in NSTX-U,'' \textit{32nd Symposium on Fusion Technology (SOFT)}, Dubrovnik, Croatia, September 18-23, 2022.}

\CVItem{[20]}{H. Al Khawaldeh, B. Leard, S.T.  Paruchuri, T. Rafiq, E. Schuster, ``Model-based LQI Controller for Simultaneous Regulation of the Current Profile and Normalized Beta in NSTX-U,'' \textit{11th ITER International School on ITER Plasma Scenarios and Control}, San Diego, CA, USA, July 25-29, 2022.}

\CVItem{[21]}{H. Al Khawaldeh, A. Pajares, T. Rafiq, E. Schuster, ``COTSIM-based Optimization of Transport Diffusivity Model for Predictive Simulations of NSTX-U Discharges,'' \textit{63rd Division of Plasma Physics (DPP) Annual Meeting of the American Physical Society (APS)}, Pittsburgh, PA, USA, November 8-12, 2021. }

\CVItem{[22]}{Y. Manjetta, H. Al Khawaldeh, T. Rafiq, E. Schuster, ``COTSIM-based Gain Optimizer for Feedback Control of Advanced Scenarios in NSTX-U,'' \textit{63rd Division of Plasma Physics (DPP) Annual Meeting of the American Physical Society (APS)}, Pittsburgh, PA, USA (Remote), November 8-12, 2021 }

% ---------------------------------------------------------
\CVSection{Research Activities}
\CVSubtitle{Collaborative Experimental Campaigns at Major Fusion Facilities}
\CVItem{\vspace{-2.75mm}March 2026}{\textbf{Advanced Density Regulation for ITER and Future Reactors by Simultaneous Gas Puffing and Pellet Injection}, KSTAR (Korea Superconducting Tokamak Advanced Research), Korea Institute of Fusion Energy, Daejeon, South Korea
  \begin{itemize}
      \item Implemented an adaptive control framework in C within the real-time Plasma Control System, satisfying a ${\sim}50~\mu s$ cycle-time constraint while modulating control effort to compensate for actuator input delays of ${\sim}50$--$200~ms$.
      \item Synchronized discrete and continuous actuators within a unified control framework by designing an integral-flow logic block to infer pellet injection state, enabling coordinated closed-loop gas and pellet fueling.
      \item Validated closed-loop density regulation across 8 experimental runs on KSTAR, achieving sub-$5\%$ steady-state regulation error in collaboration with the KFE team.
  \end{itemize}
  }
\CVItem{\vspace{-2.75mm}March 2026}{%
  \textbf{Proximity Control and Safe Reinforcement Learning Algorithms for Plasma Regulation with Stability-Limit Avoidance}, KSTAR (Korea Superconducting Tokamak Advanced Research), Korea Institute of Fusion Energy, Daejeon, South Korea
  \begin{itemize}
    \item Implemented a real-time RL controller in C to enable safe closed-loop exploration of high-density plasma scenarios.
    \item Distributed the RL policy and real-time control loop across two CPUs with inter-processor communication, coupling a ${\sim}500~\mu s$ RL inference subroutine to a ${\sim}50~\mu s$ control cycle to prevent control dropouts and system crashes.
    \item Demonstrated zero safety violations across 12 closed-loop experimental runs while tracking a time-varying density target, validating safety-prioritized control.
  \end{itemize}
}
\CVItem{\vspace{-2.75mm}June 2025}{\textbf{Advanced Density Control in DIII-D for ITER and Beyond}, DIII-D National Fusion Facility, General Atomics, San Diego, CA, USA
  \begin{itemize}
     \item Prepared control-oriented models, simulation studies, and experimental procedures to support deployment of a model-reference adaptive density controller on DIII-D.
     \item Coded and validated the adaptive controller in C within the real-time Plasma Control System, satisfying a ${\sim}50~\mu s$ cycle-time constraint across multi-actuator coordination logic and actuator delay compensation.
     \item Evaluated closed-loop control performance across multiple experimental scenarios in collaboration with the GA team.
  \end{itemize}
}
% ---------------------------------------------------------
\CVSection{Teaching and Supervision}

\CVItem{\hspace{-5mm}\mbox{2021--Ongoing}}{%
  \textbf{Student Mentor} 
  \begin{itemize}
    \item Supported junior students in developing simulation frameworks using physics-based and neural-network surrogate models for modeling and predictive analysis.
    \item Mentored junior students in the design and implementation of model-based and data-driven control strategies.
  \end{itemize}
}
\CVItem{2020--2021}{\textbf{Peer Tutor} \begin{itemize}
  \item Provided one-on-one and small-group tutoring in core mechanical engineering courses, supporting students with fundamental concepts and problem-solving techniques.
  \item Assisted students by reviewing coursework, clarifying challenging topics, and guiding them through analytical and computational problems.
  \item Developed study guides and shared effective study strategies to reinforce foundational engineering knowledge and improve academic performance.
\end{itemize}}

% ---------------------------------------------------------
\CVSection{Academic and Industry Engagements}
\CVItem{\hspace{-5mm}\mbox{2020--Ongoing}}{\textbf{Research Assistant, Plasma Control Group, Lehigh University, Bethlehem, PA}
\begin{itemize}
  \item Engineered large-scale data pipelines from 537 high-fidelity simulations across 41 experiments --- totaling 286{,}477 time slices with 1{,}147 variables per slice --- to train deep neural network surrogate models for real-time control systems.
  \item Designed and trained 5 deep neural network surrogate models in Python to approximate high-fidelity physics calculations, achieving $\sim\!10^2$--$10^4\times$ faster predictions ($R^2 = 0.95$, $RMSE = 0.22$) --- enabling real-time deployment in compute-constrained environments.
  \item Constructed physics-informed models of plasma dynamics --- combining first-principles and data-driven models to predict system behavior  --- enabling $\sim\!10^2$--$10^3\times$ faster control-architecture testing and a sim-to-real transfer pipeline.
  \item Designed and benchmarked 4 multi-actuator control architectures (MPC, adaptive, RL-based, LQI) with Kalman-based estimation to identify deployment-ready solutions for accessing and sustaining high-performance plasma scenarios under safety constraints.
  \item Integrated and deployed adaptive and safe RL controllers into 2 live experimental control systems alongside engineering and physics teams at General Atomics (GA), Korea Institute of Fusion Energy (KFE), and Princeton Plasma Physics Laboratory (PPPL).
  \item Mentored 3 junior researchers on surrogate-model and control-development workflows, covering large-dataset preparation, deep neural network training, and simulation integration --- all delivered findings at peer-reviewed conferences.
\end{itemize}
}
\CVItem{\hspace{-5mm}\mbox{June--Aug 2020}}{\textbf{Research Assistant, Data for Impact, Lehigh University, Bethlehem, PA} \begin{itemize}
  \item Investigated the mechanical behavior of von Willebrand factor under applied forces to better understand its role in bleeding disorders.
  \item Performed computational simulations to study protein deformation and elongation under varying mechanical conditions.
  \item Analyzed simulation results to identify and generalize relationships between molecular-scale mechanics and clinically relevant bleeding behavior.
\end{itemize}}
\CVItem{\hspace{-5mm}\mbox{Jan--Dec 2020}}{\textbf{Consultant, Capstone: Dorman Industries, Lehigh University, Bethlehem, PA} \begin{itemize}
  \item Led a multidisciplinary student team to design an improved diesel engine exhaust gas recirculation (EGR) cooler using CAD tools and heat transfer principles.
   \item Developed and iterated detailed CAD models for a redesigned EGR cooler to address thermal and structural performance limitations.
 \item Conducted finite element (FEA) and computational fluid dynamics (CFD) analysis, resulting in a 9.77\% increase in structural yield strength and a 23.8\% reduction in noxious emissions.
  \item Collaborated with process engineers to assess manufacturing feasibility and cost implications of implementing the proposed design.
\end{itemize}}

% ---------------------------------------------------------
\CVSection{Awards}

\CVItem{\hspace{-5mm}\mbox{2021}}{\textbf{Admitted into Lehigh President's Scholar Program, which offers students with academic excellence one year of tuition-free study.}}
\CVItem{\hspace{-5mm}\mbox{2017}}{\textbf{Received the King's Scholar Scholarship, which is a full-ride award covering the total cost of attendance for a four-year undergraduate program.}}
% % ---------------------------------------------------------
% \CVSection{Hobbies}

% \CVItem{\hspace{-5mm}\mbox{2025--Ongoing}}{\textbf{Underwater Robotics, Lehigh University, Bethlehem, PA} \begin{itemize}
%   \item Helped a group of undergraduate in designing and assembling a suitable structure for an underwater robot.
% \end{itemize}}
% ---------------------------------------------------------
% \CVSection{References}
% \CVItem{}{\textbf{Available on Request.}}
\end{document}